% Einheit 2 von 4 – Winkel zwischen Vektoren, Kanten und Geraden (bank/skalarprodukt-und-winkel/e2.jsonl, zusammenbau v0.1)
%% TODO Zeitmarke Einheit 2: Marken-Zeile nicht lesbar
%% TODO Zweigzeile Teil 1: was hier gelernt wird (Satz in Schülersprache aus dem Einheitstitel) – Einheit 2
\einheitenkopf[e2]{Einheit 2 von 4 · Winkel zwischen Vektoren, Kanten und Geraden}
\zweigzeile{Abitur GK}
\setcounter{aufgabe}{10}

%% TODO Titel: Ich-kann-Satz fehlt in der Bank (Platzhalter: Kettenname) – Nr. 11
%% TODO Anweisung: ein Satz über der ersten Teilaufgabe fehlt in der Bank – Nr. 11
\begin{aufgabe}{Der Formelwinkel oder sein Nachbar?}
\begin{teile}
\teil Der Winkel, den die Rechnung liefert, heißt Formelwinkel. Für den Innenwinkel des Dreiecks PQR bei P werden $\overrightarrow{PQ}$ und $\overrightarrow{PR}$ in die Kosinusformel eingesetzt. Welcher Winkel ist die Antwort? \\ \kreuz{Formelwinkel} \\ \kreuz{Nebenwinkel} \\ \kreuz{Komplementwinkel}
\end{teile}
\end{aufgabe}

%% TODO Titel: Ich-kann-Satz fehlt in der Bank (Platzhalter: Kettenname) – Nr. 12
%% TODO Anweisung: ein Satz über der ersten Teilaufgabe fehlt in der Bank – Nr. 12
\begin{aufgabe}{Winkel zwischen Vektoren, Kanten und Geraden}
%% TODO \rechenplatz{n}: Schrittzahl des Grundfalls fehlt in der Bank (nur bei fester Schreibform, 2.3 b)
\begin{teile}
\teil Gesucht ist der Innenwinkel des Dreiecks ABC bei A. Welche zwei Vektoren setzt du an? \\ \kreuz{$\overrightarrow{AB}$ und $\overrightarrow{AC}$} \\ \kreuz{$\overrightarrow{BA}$ und $\overrightarrow{AC}$} \\ \kreuz{$\overrightarrow{AB}$ und $\overrightarrow{BC}$}
\teil Gegeben sind $P(1 \mid 2 \mid 0)$, $Q(3 \mid 3 \mid 2)$ und $R(4 \mid 0 \mid 1)$. Berechne den Winkel zwischen den Kanten PQ und PR. $\varphi \approx$ \leerfeld[$^\circ$]
\teil Das Viereck ABCD mit $A(1 \mid 0 \mid 0)$, $B(5 \mid 2 \mid 0)$, $C(4 \mid 4 \mid 2)$ und $D(2 \mid 3 \mid 2)$. ist ein Trapez mit $AB \parallel DC$. Berechne den Innenwinkel bei D. $\varphi \approx$ \leerfeld[$^\circ$]
\teil Eine dreieckige Dachfläche hat die Ecken $P(2 \mid 0 \mid 5)$, $Q(6 \mid 2 \mid 4)$ und $R(3 \mid 5 \mid 6)$. Berechne den Innenwinkel des Dreiecks PQR bei Q. $\varphi \approx$ \leerfeld[$^\circ$]
\teil Gegeben sind $A(2 \mid 1 \mid 0)$, $B(4 \mid 3 \mid 1)$ und $C(5 \mid 1 \mid 3)$. Weise nach, dass das Dreieck ABC rechtwinklig ist, und gib die Ecke mit dem rechten Winkel an.
\teil Die Geraden $g\colon \vec x = (1 \mid 0 \mid 2) + t \cdot (2 \mid -1 \mid 2)$ und $h\colon \vec x = (3 \mid -1 \mid 4) + s \cdot (-2 \mid 2 \mid 1)$ mit $t, s \in \mathbb{R}$ schneiden sich. Berechne ihren Schnittwinkel. $\varphi \approx$ \leerfeld[$^\circ$]
\teil Das Dreieck ABC hat die Ecken $A(1 \mid 2 \mid 0)$, $B(3 \mid 4 \mid 0)$ und $C(3 \mid 2 \mid 2)$; sein Innenwinkel bei A beträgt $60^\circ$. Bestimme die beiden anderen Innenwinkel. \hfill (Abitur 2023 LK)
\end{teile}
\end{aufgabe}

%% TODO Titel: Ich-kann-Satz fehlt in der Bank (Platzhalter: Kettenname) – Nr. 13
%% TODO Anweisung: ein Satz über der ersten Teilaufgabe fehlt in der Bank – Nr. 13
\begin{aufgabe}{Winkel zwischen Vektoren, Kanten und Geraden – Fehler finden, Begründen, Anwendung}
\begin{teile}
\teil Gesucht ist der Innenwinkel des Dreiecks PQR bei P mit $P(1 \mid 0 \mid 2)$, $Q(3 \mid 2 \mid 3)$ und $R(2 \mid 0 \mid 6)$. Kai rechnet: $\overrightarrow{QP} = (-2 \mid -2 \mid -1)$, $\overrightarrow{PR} = (1 \mid 0 \mid 4)$, $\mathrm{cos}\,\varphi = \frac{-6}{3 \cdot \sqrt{17}}$, $\varphi \approx 119{,}0^\circ$. Finde den Fehler.
\teil Begründe, warum für den Innenwinkel eines Dreiecks bei A beide Vektoren von A weg zeigen müssen.
\teil Ein Fußballtor ist $7{,}32$ m breit, seine Pfosten stehen in $P(0 \mid -3{,}66 \mid 0)$ und $Q(0 \mid 3{,}66 \mid 0)$; 1 LE entspricht 1 m. Eine Spielerin schießt vom Punkt $S(16 \mid 8 \mid 0)$ aus. Unter welchem Winkel sieht sie das Tor? $\varphi \approx$ \leerfeld[$^\circ$]
\end{teile}
\end{aufgabe}
